% PRACTICA No. 2: MANEJO Y USO DEL DOSÍMETRO DE LECTURA DIRECTA

\textbf{OBJETIVO:}
Al término de esta práctica, el alumno debe reconocer que los dosímetros de lectura directa resultan ser una herramienta de seguridad radiológica de gran utilidad para todo radiógrafo.

\textbf{MATERIAL:}
\begin{itemize}
    \item Dosímetro de lectura directa (por alumno).
    \item Cargador de dosímetro de lectura directa.
    \item Contenedor de Gammagrafía.
    \begin{itemize}
        \item[\textbullet] Conteniendo una fuente radiactiva con la suficiente actividad como para poder tener niveles de rapidez de exposición a contacto de alrededor de 50 a 100 mR/h.
    \end{itemize}
    \item Cronómetro.
\end{itemize}

\textbf{DESARROLLO:}
Se le debe entregar a cada alumno un dosímetro de lectura directa, el cual tendrá que revisar adecuadamente verificando que su carga sea correcta; si no es así, deberá cargarlo utilizando para ello el cargador que anticipadamente se le mostró, hasta lograr que marque ceros. Una vez que haya logrado esto, colocará éste a un lado o encima del contenedor con la fuente radiactiva y en un campo de rapidez de exposición conocida.

Se colocará el mayor número de dosímetros posibles, controlando el tiempo de permanencia de cada uno de ellos junto al contenedor en cuestión; se irá retirando un dosímetro cada 10 minutos, anotando la lectura registrada por éste (considerando el factor de calibración del dosímetro). Con los valores obtenidos se hará una gráfica que debe coincidir con la gráfica que se haga de dosis contra tiempo, con los datos de la rapidez de exposición leídos a un costado del contenedor. Con estos datos el técnico formulará sus conclusiones con respecto al beneficio que pueda darle el adecuado uso de estos equipos.